% 【本文件写什么】子模块 runtime-heap-allocator，\subsection 级聚合
% - 职责摘要：内核堆：buddy，默认/ TLSF。
% - 子域聚合 lib.rs：os/components/wateros-runtime/runtime-heap-allocator/src/lib.rs
% - 如何重导出 api-v0、feature 下挂载哪些 impl
% - 被谁依赖，syscall、vfs、main 启动链等
% - 短综述 + 下方 \input 展开 api 与 impl 叶子
% 【事实来源】
% - os/components/wateros-runtime/runtime-heap-allocator/
% - docs/exports/ 中与 wateros-runtime 相关的 public-api / features 章节

\subsection{runtime-heap-allocator}

\code{wateros-runtime-heap-allocator}实现内核全局\code{GlobalAlloc}：默认 \code{impl-linked-list-allocator}，即 \code{linked\_list\_allocator::LockedHeap}，可选 \code{impl-tlsf}，即 \code{rlsf::Tlsf}，二者互斥由\code{compile\_error!}守卫。堆池大小\code{KERNEL\_HEAP\_SIZE}来自\code{wateros-base-config}；静态\code{HEAP\_SPACE}数组可由链接脚本映射到\code{kernel\_heap}符号。

\code{init()}在单核引导路径、堆尚未使用时初始化后端区域。\code{heap\_mem\_stats}返回\code{used}/\code{free}/\code{capacity}：链表后端为精确记账，TLSF后端\code{used}为估算值。\code{handle\_alloc\_error}由根\code{\#[alloc\_error\_handler]}委托，先\code{warn}打印布局与用量再\code{panic}。所有分配路径经\code{interrupt\_guard}关中断并检测递归分配；用量超过容量90\%时\code{maybe\_warn\_high\_water}发出警告。

\code{heap\_fragmentation\_stress\_report}为开发诊断入口，压测结束后\code{loop \{\}}挂起，不用于正常启动链。

\begin{rustcode}
pub struct HeapMemStats {
    pub used: usize,
    pub free: usize,
    pub capacity: usize,
}

pub fn handle_alloc_error(layout: core::alloc::Layout) -> ! {
    let stats = heap_mem_stats();
    log::warn!("[heap] OOM: layout_size={} ...", layout.size());
    panic!("Heap allocation error, layout = {:?}", layout);
}
\end{rustcode}
